% =====================================================================
%  7 Änderungen gegenüber Portfolioteil 1  (ab Phase 2)
%  Feedback der Tutor:in zu Phase 1 und Änderungen aus der Erarbeitung.
% =====================================================================
\section{Änderungen gegenüber den früheren Portfolioteilen}\label{sec:aenderungen}
Dieses Kapitel hält fest, was sich nach den Abgaben der ersten beiden Portfolioteile geändert hat und warum. Die meisten Änderungen gehen auf das Feedback der Tutor:in zurück, die übrigen auf Erkenntnisse aus Implementierung, Review und Tests.

\subsection{Änderungen nach Portfolioteil 1}
Tabelle~\ref{tab:aenderungen} listet die Änderungen seit der Abgabe von Portfolioteil~1. Die ersten sechs Zeilen gehen auf das Feedback zu Phase~1 zurück, die übrigen auf Erkenntnisse aus der Erarbeitungsphase.

\begingroup
\setstretch{1.0}
\fontsize{10pt}{12pt}\selectfont
\begin{longtable}{@{}L{4.0cm}L{7.4cm}L{4.6cm}@{}}
  \caption{Änderungen nach Portfolioteil 1}\label{tab:aenderungen}\\
  \toprule
  \textbf{Feedbackpunkt bzw. Anlass} & \textbf{Änderung} & \textbf{Dokument und Abschnitt} \\
  \midrule
  \endfirsthead
  \toprule
  \textbf{Feedbackpunkt bzw. Anlass} & \textbf{Änderung} & \textbf{Dokument und Abschnitt} \\
  \midrule
  \endhead
  \bottomrule
  \endfoot
  Feedback: Projektfrage nicht ausdrücklich formuliert & Projektfrage als hervorgehobener Absatz ergänzt; dazu ein Absatz, welches Dokument welchen Teil der Frage beantwortet & Projektdokumentation, Zielsetzung und Kundennutzen \\
  Feedback: nichtfunktionale Anforderungen schwer messbar & Prüfkriterien von NFA01--NFA11 mit Zahlen versehen (Testpersonen, Abdeckung, Abbruchversuche, Spaltenzahl, Regelzahl) & Anforderungsdokument, Nichtfunktionale Anforderungen \\
  Feedback: Skalierbarkeit fehlt & neue Anforderung NFA12 mit Grenzwerten für 30.000 Aktivitäten; Messwerte im Architekturdokument & Anforderungsdokument, Nichtfunktionale Anforderungen; Architekturdokument, Schichten \\
  Feedback: Wiederholungen zwischen den Dokumenten & Produktvision auf drei Sätze mit Verweis gekürzt; Problemstellung steht nur noch in der Projektdokumentation & Anforderungsdokument, Einleitung \\
  Feedback: Theorie und Entscheidung vermischt & Theorie mit Quellen und eigene Entscheidung in getrennten Absätzen, die Theorie jeweils zuerst & Projektdokumentation, Vorgehensmodell; Spezifikationsdokument, Regelauswertung; Architekturdokument, Architekturstil \\
  Feedback: Sprache straffen & geänderte Abschnitte sprachlich gekürzt, lange Sätze geteilt & alle Dokumente \\
  Widerspruch im Datenmodell & Zeichengrenze der Antworten in der Wochenreflexion einheitlich auf 1.000 Zeichen gesetzt & Spezifikationsdokument, Datenmodell \\
  Erkenntnis aus der Implementierung & Attribut \texttt{bewegungsart} in der Wiederholungsserie, damit Serientermine für R01 und R02 zählen; JSON-Beispiel ergänzt & Spezifikationsdokument, Datenmodell und Systemschnittstellen \\
  Erkenntnis aus der Implementierung & Zeile \enquote{ohne Uhrzeit} in der Wochenansicht, Kontrollkästchen zum Bearbeiten ganzer Serien, Sprung zur Vorwoche nach dem Laden der Demodaten & Architekturdokument, Änderungen gegenüber der Konzeptionsphase \\
  Erkenntnis aus dem Review & GR18 und GR20 präzisiert: Löschen oder Verschieben eines Serientermins legt eine Ausnahme an, einzeln bearbeitete Termine bleiben beim Löschen der Serie erhalten & Spezifikationsdokument, Geschäftsregeln \\
  Einsatz von \ac{KI}-Werkzeugen & Hinweis zu Umfang und Regeln des Einsatzes von Claude Code & Projektdokumentation, Vorgehensmodell und Werkzeuge \\
  Fortschritt der Phase~2 & Kapitel Versionsverwaltung und Fortschrittskontrolle zu M2 ergänzt & Projektdokumentation, Versionsverwaltung und Zeitplanung \\
\end{longtable}
\endgroup
\quelle{Eigene Darstellung.}

\abPhase{3}{%
\subsection{Änderungen nach Portfolioteil 2}
Das Feedback zu Phase~2 bewertete Struktur, Methodik und Formalien positiv und nannte vier Punkte für die Finalisierung. Tabelle~\ref{tab:aenderungen-p2} zeigt, wie sie umgesetzt wurden, und ergänzt die Änderungen aus der Testphase.

\begingroup
\setstretch{1.0}
\fontsize{10pt}{12pt}\selectfont
\begin{longtable}{@{}L{4.0cm}L{7.4cm}L{4.6cm}@{}}
  \caption{Änderungen nach Portfolioteil 2}\label{tab:aenderungen-p2}\\
  \toprule
  \textbf{Feedbackpunkt bzw. Anlass} & \textbf{Änderung} & \textbf{Dokument und Abschnitt} \\
  \midrule
  \endfirsthead
  \toprule
  \textbf{Feedbackpunkt bzw. Anlass} & \textbf{Änderung} & \textbf{Dokument und Abschnitt} \\
  \midrule
  \endhead
  \bottomrule
  \endfoot
  Feedback: Limitationen, kritische Bewertung und Fazit fehlen & neues Kapitel mit Beantwortung der Projektfrage, Limitationen a posteriori, kritischer Bewertung des Vorgehens sowie Fazit und Ausblick & Projektdokumentation, Kapitel~8 \\
  Feedback: nichtfunktionale Anforderungen stärker operationalisieren & jedes Prüfkriterium nennt Messgröße, Grenzwert und Prüfverfahren mit Testfall; NFA01 prüft statt Testpersonen automatisiert messbare Merkmale der Oberfläche & Anforderungsdokument, Kapitel~4; Spezifikationsdokument, Kapitel~5 \\
  Feedback: Wiederholungen zwischen den Dokumenten & Gesundheitsrichtwerte, Persona, Could- und Won't-Listen und doppelte Definitionen durch Verweise auf die Stelle ersetzt, an der sie ausführlich stehen & Projektdokumentation, Kapitel~1 und~6; Spezifikationsdokument, Kapitel~2; Anforderungs- und Architekturdokument \\
  Feedback: Prägnanz einzelner Formulierungen & lange Sätze geteilt, Füllwörter gestrichen, Verweise präzisiert & alle Dokumente \\
  Erkenntnis aus den Systemtests & Eingabefelder ohne verknüpfte Beschriftung ergänzt; Rückfragen zeigen deutsche Knöpfe statt \enquote{Yes} und \enquote{No} & Anwendung; Testdokument, Kapitel~3 \\
  Auslieferung & ausführbare Datei als einzelne Datei ohne ungenutzte Qt-Module in \texttt{app/release} & Projektdokumentation, Abschnitt~5.1 \\
  Fortschritt der Phase~3 & Testdokument, Benutzeranleitung, Anforderungserfüllung und Soll-Ist-Vergleich des Zeitplans ergänzt & Testdokument; Projektdokumentation, Kapitel 3, 5 und 6 \\
\end{longtable}
\endgroup
\quelle{Eigene Darstellung.}
}

